\documentclass[12pt]{article}
\usepackage[utf8]{inputenc}
\usepackage[T1]{fontenc}
\usepackage[margin=1in]{geometry}
\usepackage{setspace}
\usepackage{xcolor}
\usepackage[colorlinks=true, urlcolor=blue!60!black, linkcolor=blue!60!black]{hyperref}
\setlength{\parindent}{2em}
\setlength{\parskip}{1.25em}
\renewcommand{\baselinestretch}{1.0}
\begin{document}
\begin{center}
{\LARGE \textit{National Republican Senatorial Committee v. Federal Election Commission}, 609 U.S. \_\_\_ (2026)\par}
\vspace{0.8em}
{\large Case Brief --- Four-Page Edition\par}
\end{center}
\vspace{1em}

\section{Memory Jogger}

Overruling \textit{Colorado II}, the Court held 6--3 that FECA's limits on political parties' coordinated expenditures with their own candidates violate the First Amendment: base contribution limits, earmarking rules, and disclosure already address \textit{quid pro quo} corruption, making the caps a disproportionate ``fourth prophylaxis'' on core political speech (\href{https://www.supremecourt.gov/opinions/25pdf/24-621_h315.pdf}{Slip Opinion}).

\section{Detailed Case Facts}

FECA, 52 U.S.C. \S~30116(d), caps what a party committee may spend in ``coordination'' with its own federal candidates. As of 2026, a national party could spend between \$130,600 and \$4,071,800 in coordination with a Senate candidate, and between \$65,300 and \$130,600 with a House candidate. \textit{Colorado I} (1996) let parties make unlimited \textit{independent} expenditures, but \textit{Colorado II} (2001) upheld the coordinated caps 5--4 as closely drawn to prevent donors from circumventing candidate contribution limits through parties. After \textit{Citizens United} and \textit{SpeechNow.org}, Super PACs could raise and spend unlimited sums; in the 2023--2024 cycle, PACs raised over \$15.7 billion compared to \$2.7 billion raised by parties. In 2022, the NRSC, the NRCC, then-Senate candidate JD Vance, and then-Representative Steve Chabot sued the FEC. Vance, later Vice President, kept an active Statement of Candidacy for a 2028 Senate run.

\section{Procedural History}

The Southern District of Ohio certified the constitutional question to the en banc Sixth Circuit, which rejected the challenge in 2024. Chief Judge Sutton held the court bound by \textit{Colorado II}, though a majority doubted its vitality after \textit{McCutcheon} and \textit{Cruz}; Judge Readler dissented. In May 2025 the federal government declined to defend the law. The Supreme Court granted certiorari, allowed the Democratic committees to intervene, and appointed Roman Martinez as \textit{amicus} to defend the judgment. It decided the case on June 30, 2026.

\section{Judicial Votes}

Reversed and remanded, \textbf{6--3}.
\begin{itemize}
\item \textbf{Majority:} Justice Kavanaugh, joined by Chief Justice Roberts and Justices Thomas, Alito, Gorsuch, and Barrett.
\item \textbf{Concurrences:} None.
\item \textbf{Dissent:} Justice Kagan, joined by Justices Sotomayor and Jackson.
\end{itemize}

\section{Holding}

Section 30116(d)'s party coordinated-expenditure limits violate the First Amendment, and \textit{Colorado II} is overruled. Even under ``closely drawn'' scrutiny, a restriction must serve the interest in preventing \textit{quid pro quo} corruption or its appearance and be necessary, narrowly tailored, and not disproportionate. The case was not moot, because FECA authorizes private enforcement suits and Vance remained a declared candidate. The Court called the choice of scrutiny ``academic,'' and footnote 6 reserved limits on coordinated expenditures by outside groups.

\section{Analysis of Opinions}

\textbf{Majority (Kavanaugh).} Mass political communication requires money, so limiting a party's coordinated spending abridges its speech and has a ``stifling effect on the ability of the party to do what it exists to do.'' For nearly two centuries parties coordinated freely with their nominees. Of four asserted interests, reducing spending and curbing a party's or donors' ``undue influence'' are impermissible under \textit{Buckley}, \textit{McCutcheon}, and \textit{Cruz}. As to circumvention, a gift to a party is a step removed from corruption because the party controls the money; earmarked gifts already count against the candidate limit; disclosure exposes large gifts; and most States impose no caps without evident corruption. \textit{Colorado II} was a ``three-legged stool where all three legs have already been knocked out'': egregiously wrong, harmful in weakening parties relative to Super PACs, and supported by no cognizable reliance. Existing earmarking rules answer the joint-fundraising concern.

\textbf{Kagan (dissenting).} Both sides accept that base limits prevent corruption and that Congress may prevent their evasion. In practice candidates forward their bills to the party, making coordinated spending ``as useful to the candidate as cash.'' Through a joint fundraising committee with the national and 50 state parties, one donor can write a \$551,300 check, nearly 80 times the \$7,000 limit, and the party can spend it all on the soliciting candidate. Unearmarked checks escape the earmarking rule, and disclosure cannot deter corruption in plain sight. Nor was there a ``special justification'' for overruling: any weakening of parties stems from \textit{Citizens United} and \textit{SpeechNow.org}.

\section{Examples: Future Applications}

\begin{enumerate}
\item \textbf{Same-Side (Invalid): State Party Caps.} A State caps its parties' coordinated spending with a gubernatorial nominee while also enforcing base limits, earmarking, and disclosure. The cap is a disproportionate ``fourth prophylaxis.''
\item \textbf{Same-Side (Invalid): Digital and Field Caps.} Congress caps only coordinated polling, micro-targeting, and get-out-the-vote work. The majority treated such work as traditional party-candidate coordination.
\item \textbf{Opposite-Side (Upheld): Super PAC Coordination.} A Super PAC challenges the limit on outside groups' coordinated expenditures. Footnote 6 reserved that question, and parties' ``unity of interest'' with nominees is unique.
\item \textbf{Opposite-Side (Upheld): Earmarked Party Gifts.} The FEC fines a donor who sent \$250,000 to a party ``strictly to pay for Senator X's October TV buy.'' Earmarking enforcement remains valid.
\item \textbf{Fence-Sitter: Direct Party Cash Transfers.} A party challenges the \$5,000 cap on direct cash transfers to its candidate. Unlimited bill-paying makes the cap look anomalous, but direct cash has no expressive content and poses a more direct conduit risk.
\end{enumerate}

\section{Critique}

\textbf{Support across camps.} Originalists praised fidelity to the First Amendment's text and the long tradition of unrestricted party coordination before 1974. Even reform-minded scholars such as Samuel Issacharoff, Richard Pildes, and Bob Bauer welcomed the result; Issacharoff called it ``a welcome and overdue correction'' that rebalances campaign finance away from unaccountable outside groups.

\textbf{Progressive and doctrinal critique.} The Brennan Center and the Constitutional Accountability Center argue that Congress, not the Court, should make these structural trade-offs and that the Court again brushed aside \textit{stare decisis}. Echoing Kagan, they warn that combining \textit{McCutcheon} with unlimited coordinated spending lifts a donor's effective ceiling from \$7,000 to over \$550,000.

\section{Key Quotations}

\begin{quote}
``To sum up: In light of the other meaningful prophylactic measures available to the Government, and given the severe infringement on First Amendment-protected political speech that ensues from limiting a political party's spending in support of its candidates, we conclude that the political-party coordinated-expenditure limits are `disproportionate' and are not `necessary' and `narrowly tailored' for the circumvention interest it seeks to protect.''
\end{quote}
\noindent --- Justice Kavanaugh, for the Court, slip op. at 21.

\begin{quote}
``The base limits on contributions to candidates serve as an initial prophylaxis against quid pro quo corruption or its appearance in this context---after all, most contributions to candidates are not given in exchange for some official action. \textit{Id.}, at 221. The earmarking rules constitute a second prophylaxis. The disclosure requirements supply a third prophylaxis. So prophylaxis upon prophylaxis upon prophylaxis already serve to prevent quid pro quo corruption or its appearance.''
\end{quote}
\noindent --- Justice Kavanaugh, for the Court, slip op. at 20.

\begin{quote}
``The majority invalidates Congress's restriction of coordinated expenditures, thus enabling a party to serve as an alternative checking account for a campaign. As a result, a donor will be able to give a party as much as half a million dollars (as compared to the \$7,000 he can give directly to the candidate) to cover the candidate's bills.''
\end{quote}
\noindent --- Justice Kagan, dissenting, slip op. at 2.

\end{document}
